\begin{helper}
In the fixed exposure-history cell of
\(\noderef{FixedSetHistoryCellFixedCountProductBound}\), put
\[
a=f(\ell_R,\ell_B)-10\varepsilon_1 k^2 .
\]
For fixed nonnegative counts \(u_R,u_B\), the finite branch/transcript
decomposition gives the following two ordered product bounds.  If
\(u_B\le u_R\), then there are natural numbers \(N_R,N_B\) with
\[
N_R\le 3\varepsilon_1 k^2,\qquad N_B\le k^2
\]
such that the conditional probability of the fixed-count target is at most
\[
\binom{N_R}{u_R}p^{u_R}\binom{N_B}{u_B}p^{u_B}
(1-p)^{\max(0,a-u_R-u_B)}.
\]
If \(u_R\le u_B\), then the same statement holds after interchanging the
colors:
\[
\binom{N_R}{u_B}p^{u_B}\binom{N_B}{u_R}p^{u_R}
(1-p)^{\max(0,a-u_B-u_R)}
\]
with \(N_R\le3\varepsilon_1 k^2\) and \(N_B\le k^2\).
\end{helper}
\begin{proof}
We prove the case \(u_B\le u_R\).  The proof of the second displayed
inequality is the same argument with the two colors interchanged.

Write \(H\) for the fixed history cell and write \(T_{u_R,u_B}\) for the
event in the statement: the final \(I\)-set is independent, the regularity
event holds, and the red and blue present staged-open terminal supports have
cardinalities \(u_R\) and \(u_B\).  Put
\[
L=\max(0,a-u_R-u_B).
\]
All terminal coordinates lie in the finite set \(U\) called ``terminal'' in
the statement, and the history-cell hypotheses say that \(H\) fixes the
embedding and every recorded coordinate.  Thus, after conditioning on \(H\),
the only remaining random coordinates relevant to this argument are the
terminal coordinates, in the fixed order supplied in the statement.  The
terminal product law is exactly the product law used in
\(\noderef{FixedSetHistoryCellTerminalProductPartition}\).

First choose the finite blue support universe.  Let \(U_B\) be the set of
blue terminal projected pairs that can arise from pairs of vertices of
\(I\).  The pair set on \(I\) has at most \(k^2\) ordered pairs, so
\[
N_B:=|U_B|\le k^2.
\]
Every blue present staged-open support of a sample in \(H\cap T_{u_R,u_B}\)
is a subset \(B\subseteq U_B\) of size \(u_B\).  Hence the number of possible
blue supports is at most \(\binom{N_B}{u_B}\).

Fix such a blue support \(B\).  We next partition the samples with blue
support \(B\) by their complete blue terminal trace.  A trace \(\sigma\) is a
choice, for every blue terminal coordinate \(e\), of whether
\(X_e=1\) or \(X_e=0\).  The trace class consists of samples whose blue
terminal values agree with \(\sigma\), and whose blue present staged-open
support is exactly \(B\).  These trace classes cover the part of
\(H\cap T_{u_R,u_B}\) with blue support \(B\), and two different traces are
disjoint, by \(\noderef{FixedSetHistoryCellAdaptiveBranchPartition}\).
There are finitely many such traces because \(U\) is finite.

For every pair \((B,\sigma)\) whose trace class meets
\(H\cap T_{u_R,u_B}\), choose one representative branch
\(\beta_{B,\sigma}\) in that nonempty class.  Apply the branch exceptional-set
hypothesis to this representative.  Since the representative itself is in
\(T_{u_R,u_B}\), it is regular and agrees with itself on the blue terminal
trace; therefore the red exceptional set \(E_R(B,\sigma)\) supplied by the
hypothesis satisfies
\[
|E_R(B,\sigma)|\le 3\varepsilon_1 k^2,
\]
because \(\noderef{ClosedPairsBound}\) is the assertion
\(|E_R(B,\sigma)|\le 3\varepsilon_1 k^2\).  For any other sample in the same
trace class, the same branch hypothesis applies: it agrees with
\(\beta_{B,\sigma}\) on the embedding, recorded coordinates, and all blue
terminal coordinates, so every red present staged-open coordinate of that
sample lies in \(E_R(B,\sigma)\).

Define \(N_R\) to be the maximum of the finite set of cardinalities
\(|E_R(B,\sigma)|\), over all nonempty pairs \((B,\sigma)\) just described;
if there is no such pair, set \(N_R=0\).  The preceding bound for every
nonempty pair gives
\[
N_R\le 3\varepsilon_1 k^2.
\]
Moreover, for each fixed nonempty \((B,\sigma)\), the red support \(R\) of
any sample in the trace class is a \(u_R\)-element subset of
\(E_R(B,\sigma)\).  Since \(|E_R(B,\sigma)|\le N_R\), padding
\(E_R(B,\sigma)\) inside an \(N_R\)-element set injects its \(u_R\)-subsets
into the \(u_R\)-subsets of that \(N_R\)-element set.  Thus, uniformly in
\((B,\sigma)\), the number of possible red supports is at most
\(\binom{N_R}{u_R}\).  Fix such a padding injection for every nonempty
\((B,\sigma)\).  We use its image as a global red-support label
\(\rho\), so the final summation is indexed by \(B\) and \(\rho\), while the
complete traces \(\sigma\) are still summed through their cylinder weights.

Now fix \(B\), a nonempty trace \(\sigma\), and a red support
\[
R\subseteq E_R(B,\sigma),\qquad |R|=u_R.
\]
Inside this fixed branch and support class, apply
\(\noderef{FixedSetHistoryCellCanonicalAbsenceSelection}\) to each successful
sample.  Its hypotheses are satisfied as follows.  The regularity and
projection-size assumptions give
\[
a\le |\noderef{TwoBiteStagedOpenPairs}(\omega,\varepsilon,I)|,
\]
all staged-open coordinates are terminal by the history-cell containment
hypothesis, all present blue staged-open coordinates are in \(B\), and all
present red staged-open coordinates are in \(R\) by the exceptional-set
argument above.  Hence the canonical selection produces a set
\[
Z(\omega)\subseteq U\setminus(B\cup R)
\]
of absent staged-open terminal coordinates such that
\[
|Z(\omega)|\ge L.
\]
Each coordinate of \(Z(\omega)\) is selected at the moment it is reached in
the fixed terminal order, and the prefix-safety hypothesis says that the
decision to select it depends only on the history and the terminal values
strictly earlier in the order.

We refine once more by prefix transcripts.  A transcript \(\tau\) records:
the fixed blue trace \(\sigma\); the fixed blue support \(B\); the fixed red
support \(R\); every terminal value inspected before and while the canonical
algorithm selects the coordinates of \(Z\); and the selected absent set
\(Z_\tau\).  In particular, every terminal value fixed by the complete trace
or by the prefix transcript is recorded in \(\tau\), including terminal
values that are not members of \(B\), \(R\), or \(Z_\tau\).  For such a
transcript define two disjoint finite sets
\[
P_\tau=\{e:\tau\text{ fixes }X_e=1\},\qquad
A_\tau=\{e:\tau\text{ fixes }X_e=0\}.
\]
Thus \(B\cup R\subseteq P_\tau\), \(Z_\tau\subseteq A_\tau\), and every fixed
terminal value from \(\sigma\) and from the prefix transcript belongs to
exactly one of \(P_\tau\) and \(A_\tau\).  The cell \(G_{B,\sigma,R,\tau}\)
is the event that all these fixed terminal values occur, together with the
history \(H\).  These cells are finite in number and cover the portion of
\(H\cap T_{u_R,u_B}\) with the fixed data \(B,\sigma,R\).  Distinct complete
traces or distinct transcripts disagree on at least one terminal value or on
the stopping data of the canonical selection, so the refinement is a genuine
finite partition, up to empty cells.

The weight attached to this cell is the raw history mass multiplied by the
corresponding terminal-cylinder factor:
\[
w_{B,\sigma,R,\tau}
=\Pr(H)\,p^{|P_\tau|}(1-p)^{|A_\tau|}.
\]
This is the weight for all terminal values fixed by the trace and transcript,
not only for the mandatory support and absence coordinates.  Since
\(P_\tau\) and \(A_\tau\) are disjoint subsets of \(U\), the fixed-cylinder
bound \(\noderef{FixedSetHistoryCellFixedCylinderCharge}\), applied using the
complete terminal product law, gives
\[
\Pr(H\cap G_{B,\sigma,R,\tau})
\le w_{B,\sigma,R,\tau}.
\]
If \(\Pr(H)=0\), both sides of this raw estimate are zero or nonnegative, and
the conditional probability is zero by
\(\noderef{TwoBiteConditionalProbability}\).  Assume for the moment that
\(\Pr(H)>0\); then dividing all raw weights by \(\Pr(H)\) gives the same
estimates inside the conditional product space.

It remains to sum these cell weights without introducing a trace or
transcript counting factor.  Fix a blue support \(B\) and a global
red-support label \(\rho\).  Consider all cells
\(G_{B,\sigma,R,\tau}\) for all complete blue traces \(\sigma\), all actual
red supports \(R\subseteq E_R(B,\sigma)\) whose padding image is \(\rho\),
and all prefix transcripts \(\tau\).  If no such actual support exists for a
trace, that trace contributes no cells.  Otherwise, view the complete trace
and the prefix transcript together as a finite decision tree in the terminal
order.  At a node of the tree, the next coordinate to inspect is determined
by the history and the previously recorded terminal values.  If that
coordinate is one of the mandatory present coordinates in the actual set
\(B\cup R\), only the present child is retained and contributes a factor
\(p\).  If it is one of the selected absent coordinates in \(Z_\tau\), only
the absent child is retained and contributes a factor \(1-p\).  Every other
terminal value fixed by the blue trace or by the prefix transcript is
auxiliary: when the leaves over all \(\sigma\) and \(\tau\) with the fixed
labels \(B,\rho\) are summed, both possible children occur and their factors
add to
\[
p+(1-p)=1.
\]
Equivalently, after removing the mandatory present set \(B\cup R\) and the
selected absent set \(Z_\tau\), the remaining fixed terminal coordinates are
reindexed by subsets \(W\) of a finite residual set, and their contribution is
\[
\sum_W p^{|W|}(1-p)^{|W^c|}=1.
\]
This is a finite partition/reindexing identity, not an omission of the trace
or transcript values: those values are included in \(P_\tau\) or \(A_\tau\)
before the summation is performed.

Consequently, after summing over all complete traces and prefix transcripts
compatible with the fixed labels \(B\) and \(\rho\), all auxiliary terminal
factors sum to at most \(1\), while the mandatory factors remain.  For every
nonempty leaf, the actual mandatory red support has size \(u_R\), and the
blue support has size \(u_B\).  Thus the mandatory present factors contribute
\[
p^{|B|}p^{|R|}=p^{u_B}p^{u_R}.
\]
The selected absent set has size at least \(L\), and \(0\le p\le 1/2\) implies
\(1/2\le 1-p\le1\).  Hence increasing the exponent can only decrease the
factor, so
\[
(1-p)^{|Z_\tau|}\le (1-p)^L.
\]
Thus the conditional sum of all transcript weights for fixed \(B\) and
\(\rho\)
is at most
\[
p^{u_B}p^{u_R}(1-p)^L.
\]

Finally sum over supports.  There are at most \(\binom{N_B}{u_B}\) choices
for \(B\), and there are at most \(\binom{N_R}{u_R}\) global red-support
labels \(\rho\).  The finite cell family, the per-cell raw estimates, and the
preceding total weight bound for each pair of labels are precisely the
hypotheses of
\(\noderef{FixedSetHistoryCellBranchTranscriptSummation}\).  That lemma gives
the raw estimate
\[
\Pr(H\cap T_{u_R,u_B})
\le
\Pr(H)\binom{N_R}{u_R}p^{u_R}\binom{N_B}{u_B}p^{u_B}(1-p)^L.
\]
When \(\Pr(H)>0\), division by \(\Pr(H)\) gives the first displayed
conditional-probability bound.  When \(\Pr(H)=0\), the conditional probability
is zero and the same inequality follows from the nonnegativity of all factors.

For the case \(u_R\le u_B\), repeat the construction with the colors
interchanged.  The complete trace is now the red terminal trace, its support
universe has size \(N_B\le k^2\), and the branch hypothesis supplies blue
exceptional sets whose finite maximum is \(N_R\le3\varepsilon_1k^2\).  The
same transcript-weight summation gives
\[
\binom{N_R}{u_B}p^{u_B}\binom{N_B}{u_R}p^{u_R}
(1-p)^{\max(0,a-u_B-u_R)},
\]
which is the second displayed bound.
\end{proof}
